\section{The actual first power sharpens common-exponent support to a split progression}
\label{acs-section}

Let $a,b$ be positive coprime integers, $c=a+b$, and suppose the actual
norms satisfy
\begin{equation}\label{acs-pure}
 M=a^2+ab+b^2=R^p,\qquad
 F=a^4+3a^3b+5a^2b^2+3ab^3+b^4=Q^p,
\end{equation}
where $p\ge5$ is prime and $R,Q$ are positive integers. For a positive
integer modulus $s$, write
\[
 c_{\mathrm{bad}}(s)=\prod_{q\not\equiv1\bmod s}q^{v_q(c)},
 \qquad c_{\mathrm{good}}(s)=c/c_{\mathrm{bad}}(s),
\]
with products over primes. These are full prime-power parts, not
radicals. This section combines the common-exponent factorization with
one specific input arm of the actual Eisenstein power.

\begin{lemma}[One input arm pays for the extra excluded primes]
\label{acs-input-arm}
There is a primitive unramified $w=x+y\zeta$ of norm $R$ with
$w^p=a+b\zeta$. Set
\[
 \ell=\begin{cases}x+y,&p\equiv1\pmod6,\\x,&p\equiv5\pmod6.
 \end{cases}
\]
Then $\ell\ne0$ and
\begin{equation}\label{acs-extra-part}
 J=\prod_{\substack{q\mid c\\q\equiv1\bmod p\\q\not\equiv1\bmod6p}}
       q^{v_q(c)}\ \mid\ |\ell|,
 \qquad 3J^2\le3\ell^2\le4R.
\end{equation}
\end{lemma}
\begin{proof}
A primitive Eisenstein integer has no inert rational factor and no pair
of conjugate split prime factors. If $3\mid M$, primitiveness gives
$a\equiv b\not\equiv0\pmod3$, and direct expansion gives $v_3(M)=1$.
This is incompatible with $M=R^p$. Each remaining oriented split
exponent is divisible by $p$, because its norm exponent is. Unique
factorization gives $a+b\zeta=u w_0^p$ with $N(w_0)=R$ and
$u\in\mu_6$. The $p$th-power map on $\mu_6$ is bijective, so a unit
$p$th root of $u$ can be absorbed in $w_0$. The resulting $w$ is
primitive, unramified and nonunit; hence $R\ge7$ and its boundary
arms are nonzero.

Apply Lemma~\ref{sac-quotients} at exponent $p$. If $p\equiv1\pmod6$,
the sum of the normalized output coordinates is $c$, with input arm
$x+y$. If $p\equiv5\pmod6$, the normalized output is
$\overline{a+b\zeta}$, whose first coordinate is $c$, with input arm
$x$. Thus for one actual quotient $D_j$,
\begin{equation}\label{acs-actual-quotient}
 c=\ell D_j,\qquad T_p/T_1=|D_1D_2D_3|.
\end{equation}

Let $q$ occur in $J$. It is at least $2p+1>3$, so $q\ne p$.
If it split in the Eisenstein field, then $q\equiv1\pmod3$;
together with its oddness and $q\equiv1\pmod p$, this would give
$q\equiv1\pmod{6p}$. Hence it is inert. Also
$\gcd(M,c)=1$, so $q\nmid R$.

The actual homogeneous rank
$d_q=\operatorname{ord}((w/\bar w)^3\bmod q)$ divides $p$, since
$q\mid T_p$. In the inert case the norm-one group has order $q+1$,
and its cube subgroup has order $(q+1)/3$. Rank $p$ would force
$p\mid(q+1)/3$, impossible because $q+1\equiv2\pmod p$.
Therefore $d_q=1$. The homogeneous valuation law, as used in
Theorem~\ref{cp-valuations}, gives
\[
 v_q(T_p)=v_q(T_1)+v_q(p)=v_q(T_1).
\]
By \eqref{acs-actual-quotient}, all three nonnegative quotient
valuations are zero. Hence $v_q(c)=v_q(\ell)$, with its full depth,
and $J\mid|\ell|$. Finally
\[
 4R-3(x+y)^2=(x-y)^2,\qquad
 4R-3x^2=(x+2y)^2
\]
proves the size estimate.
\end{proof}

\begin{theorem}[Actual mass in the split progression]\label{acs-mass}
Every actual representation \eqref{acs-pure} satisfies
\begin{align}
 27c_{\mathrm{bad}}(6p)^2&\le16R^3,\label{acs-square-budget}\\
 c_{\mathrm{good}}(6p)&>
       \frac{3\sqrt3}{4}\,c^{1-3/p}.\label{acs-good-size}
\end{align}
In particular $c$ has a prime divisor $q\equiv1\pmod{6p}$, necessarily
$q\ge6p+1$, and
\begin{equation}\label{acs-weighted-mass}
 \sum_{\substack{q\mid c\\q\equiv1\bmod6p}}v_q(c)\log q
 >\left(1-\frac3p\right)\log c+\log\frac{3\sqrt3}{4}.
\end{equation}
\end{theorem}
\begin{proof}
The common-exponent allocation gives $c_{\mathrm{bad}}(p)\le2R/3$.
For completeness, put $D=Q-R^2>0$ and
$S=\sum_{j=0}^{p-1}Q^{p-1-j}R^{2j}$. Then $DS=ab c^2$.
Every prime factor of $S$ other than $p$ has exact multiplicative
order $p$ for $Q/R^2$; its only possible prime contribution outside
$1\pmod p$ is one factor of $p$. Thus the full bad part of $ab c^2$
divides $pD$. The inequalities $S\ge pR^{2p-2}$ and
$9ab c^2\le4M^2$ give $pD\le4R^2/9$, proving the asserted bound
on $c_{\mathrm{bad}}(p)$, with the exponent-prime exception included.

The supports of $c_{\mathrm{bad}}(p)$ and $J$ are disjoint, and
their product is exactly $c_{\mathrm{bad}}(6p)$. Multiplication of
$9c_{\mathrm{bad}}(p)^2\le4R^2$ and $3J^2\le4R$ proves
\eqref{acs-square-budget}. Dividing, taking square roots, and using
$R<c^{2/p}$ from $M<c^2$ proves \eqref{acs-good-size}.
Its right side exceeds one for $p\ge5$, so the good part has a
prime divisor. Logarithms give \eqref{acs-weighted-mass}.
\end{proof}

If $M=A^h$ and $F=B^g$ for positive integer bases and $h,g\ge2$,
every prime $p\ge5$ dividing $\gcd(h,g)$ gives the same conclusions
by setting $R=A^{h/p}$ and $Q=B^{g/p}$. Along a sequence with chosen
common primes $p\to\infty$, the left side of
\eqref{acs-weighted-mass} divided by $\log c$ tends to one. This is
a statement about the actual seed's valuations, with no independence
or prime-density hypothesis.

Primes already present in the input arm are fully included in its
budget. Large new primes in the surviving progression and their high
depths remain possible. The argument gives no sufficient radical
bound for ABC. Nonunit residuals, the first norm $3A^h$, and pairs
of exponents without a common prime at least five remain separate
open branches. The rank and valuation part of this section is an
ordinary proof, not a claim that the integer arm module formalizes
the entire progression argument.
